\section*{第 \ref{ch:g6:circuits-and-safety} 章——电路与安全}

\begin{solution}{exo:g6:circuits-and-safety:1}
应忠实地表示连接关系——哪个器件接到哪个器件，
以及道路在何处分叉。忽略的是导线的长度、颜色和
在桌面上蜿蜒的走向。
\end{solution}

\begin{solution}{exo:g6:circuits-and-safety:2}
\omterm{def:g3:first-electric-circuit:battery}{电池}：两条不等长的平行线；\omterm{def:g3:first-electric-circuit:bulb}{灯泡}：带十字的圆；
断开的\omterm{ex:g3:first-electric-circuit:switch}{开关}：升起的吊桥；闭合的\omterm{ex:g3:first-electric-circuit:switch}{开关}：放平的桥面；
\omterm{ex:g6:circuits-and-safety:motor}{电动机}：圆内的 M；\omterm{ex:g6:circuits-and-safety:motor}{蜂鸣器}：被敲击的圆。
较\emph{长}的那条线标出\omterm{def:g3:first-electric-circuit:battery}{电池}的 $+$ 极。
\end{solution}

\begin{solution}{exo:g6:circuits-and-safety:3}
画成一个矩形：\omterm{def:g3:first-electric-circuit:battery}{电池}、按压式\omterm{ex:g3:first-electric-circuit:switch}{开关}和\omterm{ex:g6:circuits-and-safety:motor}{蜂鸣器}位于同一回路中。
松开手指，\omterm{ex:g3:first-electric-circuit:switch}{开关}断开——唯一的通路出现缺口，
回路定律使\omterm{ex:g6:circuits-and-safety:motor}{蜂鸣器}立即停响。
\end{solution}

\begin{solution}{exo:g6:circuits-and-safety:4}
不一定。关键问题是：谁与谁相连？
器件相同，连接相同，就是同一个电路，
无论看上去的线团多么不同。
\end{solution}

\begin{solution}{exo:g6:circuits-and-safety:5}
风扇照常转动：\omterm{def:g3:first-electric-circuit:bulb}{灯泡}烧坏只使它自己的\omterm{def:g5:series-parallel-circuits:parallel}{支路}变暗——
这正是\omterm{def:g5:series-parallel-circuits:parallel}{并联}的独立性。断开的\omterm{ex:g3:first-electric-circuit:switch}{开关}位于分叉之前的公共通路上，
因此两条\omterm{def:g5:series-parallel-circuits:parallel}{支路}同时被切断：这就是总控规则。
\end{solution}

\begin{solution}{exo:g6:circuits-and-safety:6}
一根裸露的导线直接连接电源的两端，
跳过所有用电器：不受阻碍的电流猛冲过去，导线发热。
受害者：\omterm{def:g3:first-electric-circuit:battery}{电池}（耗尽并损坏）、导线的绝缘层——
在市电上，有时甚至整座房子都会遭殃。
\end{solution}

\begin{solution}{exo:g6:circuits-and-safety:7}
它在一瞬间切断了自己的\omterm{def:g5:series-parallel-circuits:parallel}{支路}，因为电流开始猛冲——
这是故障，不是心血来潮。在复位保护装置之前，
先找出并排除故障（拔掉可疑的用电器，检查该\omterm{def:g5:series-parallel-circuits:parallel}{支路}）。
\end{solution}

\begin{solution}{exo:g6:circuits-and-safety:8}
洗浴时：潮湿的皮肤导电，市电穿过湿润的身体，
可能致命。电线破损：开裂的绝缘层就像有缺口的栅栏——
手指可能碰到铜线。可以碰\omterm{def:g3:first-electric-circuit:battery}{电池}，绝不能碰插座：
\omterm{def:g3:first-electric-circuit:battery}{电池}的电压很小，插座的电压
则大数百倍。
\end{solution}

\begin{solution}{exo:g6:circuits-and-safety:9}
\omterm{def:g6:circuits-and-safety:diagram}{电路图}：\omterm{def:g3:first-electric-circuit:battery}{电池}，公共通路上一个闭合的\omterm{ex:g3:first-electric-circuit:switch}{开关}，然后分岔——
一条是\omterm{ex:g6:circuits-and-safety:motor}{蜂鸣器}\omterm{def:g5:series-parallel-circuits:parallel}{支路}，一条是\omterm{def:g3:first-electric-circuit:bulb}{灯泡}\omterm{def:g5:series-parallel-circuits:parallel}{支路}——再汇合回到 $-$ 极。预测结果：
\omterm{def:g3:first-electric-circuit:bulb}{灯泡}发光\emph{并且}\omterm{ex:g6:circuits-and-safety:motor}{蜂鸣器}鸣响，二者彼此独立。
在\omterm{ex:g6:circuits-and-safety:motor}{蜂鸣器}\omterm{def:g5:series-parallel-circuits:parallel}{支路}上加一个断开的\omterm{ex:g3:first-electric-circuit:switch}{开关}：该\omterm{def:g5:series-parallel-circuits:parallel}{支路}静音，
但\omterm{def:g3:first-electric-circuit:bulb}{灯泡}照常发光——\omterm{def:g5:series-parallel-circuits:parallel}{支路}\omterm{ex:g3:first-electric-circuit:switch}{开关}只控制自己的\omterm{def:g5:series-parallel-circuits:parallel}{支路}。
\end{solution}

\begin{solution}{exo:g6:circuits-and-safety:10}
多出的导线构成\omterm{def:g6:circuits-and-safety:short}{短路}：一条从 $+$ 到
$-$ 的裸露直通路。电流偏好这条容易走的空路，
而不去推动辛苦工作的\omterm{ex:g6:circuits-and-safety:motor}{电动机}，于是风扇缺电，
\omterm{def:g6:circuits-and-safety:short}{短路}导线却载着猛冲的电流发热——耗尽\omterm{def:g3:first-electric-circuit:battery}{电池}，还可能引起烧伤。请拆除
标着“坚固”的那根导线。
\end{solution}

\begin{solution}{exo:g6:circuits-and-safety:11}
墙上的\omterm{ex:g3:first-electric-circuit:switch}{开关}断开的是灯具自身的回路——\omterm{def:g3:first-electric-circuit:bulb}{灯泡}的通路
被切断。断路器断开的则是更靠近电源的整条\omterm{def:g5:series-parallel-circuits:parallel}{支路}——
即使灯具接线有误或墙上\omterm{ex:g3:first-electric-circuit:switch}{开关}标错，这条\omterm{def:g5:series-parallel-circuits:parallel}{支路}也会断电。
两个断口\omterm{def:g4:series-circuits:series}{串联}：单独任何一个都能切断电流，
两者同时存在时，任一处出错都有另一处兜底。
\end{solution}

\begin{solution}{exo:g6:circuits-and-safety:12}
\omterm{def:g3:first-electric-circuit:battery}{电池}，公共通路上的总\omterm{ex:g3:first-electric-circuit:switch}{开关}，然后四条\omterm{def:g5:series-parallel-circuits:parallel}{并联}\omterm{def:g5:series-parallel-circuits:parallel}{支路}：
舞台灯 1；舞台灯 2；带自有\omterm{def:g4:series-circuits:series}{串联}\omterm{ex:g3:first-electric-circuit:switch}{开关}的窗帘\omterm{ex:g6:circuits-and-safety:motor}{电动机}；
带自有\omterm{def:g4:series-circuits:series}{串联}\omterm{ex:g3:first-electric-circuit:switch}{开关}的\omterm{ex:g6:circuits-and-safety:motor}{蜂鸣器}。检验：总\omterm{ex:g3:first-electric-circuit:switch}{开关}控制全部
（公共通路）；两盏灯各自独立失效（分属不同\omterm{def:g5:series-parallel-circuits:parallel}{支路}）；
\omterm{ex:g6:circuits-and-safety:motor}{电动机}与\omterm{ex:g6:circuits-and-safety:motor}{蜂鸣器}各自只受自己的\omterm{ex:g3:first-electric-circuit:switch}{开关}控制
（\omterm{def:g5:series-parallel-circuits:parallel}{支路}内\omterm{def:g4:series-circuits:series}{串联}），彼此互不干扰
（\omterm{def:g5:series-parallel-circuits:parallel}{支路}之间\omterm{def:g5:series-parallel-circuits:parallel}{并联}）。
\end{solution}

\begin{solution}{pb:g6:circuits-and-safety:1}
\textbf{1.} \omterm{def:g4:series-circuits:series}{串联}时，一个器件损坏或被关断，
全都变暗；三个器件分担电压，个个暗弱。
\omterm{def:g5:series-parallel-circuits:parallel}{并联}则每个器件彼此独立，
且都保持亮度正常。
\textbf{2.} 每个\omterm{ex:g3:first-electric-circuit:switch}{开关}只控制其所在\omterm{def:g5:series-parallel-circuits:parallel}{支路}上的器件，
不会控制其他\omterm{def:g5:series-parallel-circuits:parallel}{支路}。
\textbf{3.} 在公共通路上，位于\omterm{def:g3:first-electric-circuit:battery}{电池}与第一个分岔之间——
位于所有\omterm{def:g5:series-parallel-circuits:parallel}{支路}的上游。
\textbf{4.} \omterm{def:g3:first-electric-circuit:battery}{电池} $\to$ 总\omterm{ex:g3:first-electric-circuit:switch}{开关} $\to$ \omterm{def:g5:series-parallel-circuits:parallel}{节点}，分出三条\omterm{def:g5:series-parallel-circuits:parallel}{支路}——
（吸顶灯 $+$ 其\omterm{ex:g3:first-electric-circuit:switch}{开关}）、（工作台灯 $+$ 其\omterm{ex:g3:first-electric-circuit:switch}{开关}）、
（风扇\omterm{ex:g6:circuits-and-safety:motor}{电动机} $+$ 其\omterm{ex:g3:first-electric-circuit:switch}{开关}）——再汇合，
回到\omterm{def:g3:first-electric-circuit:battery}{电池}。
\textbf{5.} 它们\omterm{def:g4:series-circuits:series}{串联}在同一条\omterm{def:g5:series-parallel-circuits:parallel}{支路}上，分担电压
（两者都很暗——那盏暗的吸顶灯其实是两灯\omterm{def:g4:series-circuits:series}{串联}的症状），
风扇的\omterm{ex:g3:first-electric-circuit:switch}{开关}位于它们的共用通路上，
同时控制两者。
\textbf{6.} 把它们拆开：另拉一根新导线，使工作台灯
独占一条跨接在电源两端的\omterm{def:g5:series-parallel-circuits:parallel}{支路}——两灯各自满亮，
彼此独立。
\textbf{7.} \omterm{def:g3:first-electric-circuit:battery}{接线柱}两端发生\omterm{def:g6:circuits-and-safety:short}{短路}。检查\omterm{def:g3:first-electric-circuit:battery}{电池}是否仍能供电——
\omterm{def:g6:circuits-and-safety:short}{短路}会耗尽\omterm{def:g3:first-electric-circuit:battery}{电池}，还可能损坏电芯
（发热的\omterm{def:g3:first-electric-circuit:battery}{电池}应当报废）。
\textbf{8.} \omterm{def:g6:circuits-and-safety:short}{短路}（裸导线彼此相碰）与触电、烧伤
（手指接触裸露的铜线）。
\textbf{9.} 潮湿的皮肤导电远胜干燥的皮肤；即使是
连\omterm{def:g3:first-electric-circuit:battery}{电池}这样的小电路也应养成干手操作的习惯——而在家中，
同样的接触遇上市电的电压，大数百倍：
可能致命。习惯必须在安全的电路上养成。
\textbf{10.} 断路器检测到其\omterm{def:g5:series-parallel-circuits:parallel}{支路}中电流过大——
这是故障，很可能是灯具进了雨水。正确的顺序是：让断路器保持断开，
拔掉插头并检查该\omterm{def:g5:series-parallel-circuits:parallel}{支路}，排除或烘干故障，
然后才能复位。
\textbf{11.} 它有自己的回路，独立于棚屋的电网（或作为
棚屋\omterm{def:g3:first-electric-circuit:battery}{电池}上多出的一条\omterm{def:g5:series-parallel-circuits:parallel}{并联}\omterm{def:g5:series-parallel-circuits:parallel}{支路}）：\omterm{def:g3:first-electric-circuit:battery}{电池}、住宅处的按压式\omterm{ex:g3:first-electric-circuit:switch}{开关}、
长长的双芯导线、棚屋处的\omterm{ex:g6:circuits-and-safety:motor}{蜂鸣器}，再回到\omterm{def:g3:first-electric-circuit:battery}{电池}——
门铃本身就是一个\omterm{def:g4:series-circuits:series}{串联}回路，
无论导线伸得多远。
\textbf{12.} 例如：“插座：我们的实验只到\omterm{def:g3:first-electric-circuit:battery}{电池}为止——
插座不归我们管。水：湿手不可接触任何线路，
无论室内还是室外。裸导线：每个接头都要包好；
裸导线是故障，而不是捷径。”
\end{solution}
